% SHARD-ID: AQM-75-QSA-WHAT-COUNTS-AS-A-POINT
% SHARD-TITLE: What Counts As A Point
% SHARD-SUMMARY: Compares underlying-set, rational-point, closed-point, generic-point, and all-scheme-point indexing choices in the QSA catalogue.
% SHARD-SUMMARY: Reviews the existing closed-point finite Weyl-site conventions for finite spectra, finite-type affine schemes over finite fields, and \(\Spec\mathbb Z\).
% SHARD-SUMMARY: Identifies which point choices give finite kinematical systems and records generic-point and analytic-completion gaps.
% SHARD-KEYWORDS: quantum-system association, points, underlying set, rational points, closed points, generic point, finite Weyl sites

\section{What Counts As A Point}
\label{sec:qsa-what-counts-as-a-point}

\subsection{Status and Local Anchors}

This is QSA Step 11 from
\path{docs/quantum_system_associations/PLAN.md}.  It is a review/new
synthesis shard: it adds no source, no run artifact, and no new
representation convention.  Its role is to make the point-selection choice
explicit before later geometric association shards use points as quantum
sites.

The governing comparison language is \texttt{CONVENTIONS.md} convention
\texttt{(ap)} and AQM-65.  The point conventions reviewed here are
\texttt{CONVENTIONS.md} \texttt{(t)} for the prime-field underlying-set versus
\(\Spec\mathbb F_p\) distinction, \texttt{(u)} for finite affine spectra with
residue qudits, \texttt{(y)} for closed points of affine finite-type
\(\mathbb F_q\)-schemes, and \texttt{(ai)} for \(\Spec\mathbb Z\).  The local
report anchors are AQM-22, AQM-23, AQM-28, AQM-29, AQM-30, and AQM-45, with
AQM-66 and AQM-71 used only for QSA catalogue language.

\subsection{Point-Selection Choices}

The table records choices already visible in the report.  A row is a point
convention only for the object class named in that row; it is not a universal
rule for every object called \(X\).

\begin{center}
\small
\begin{tabular}{p{0.17\linewidth}p{0.23\linewidth}p{0.24\linewidth}p{0.24\linewidth}}
\toprule
Choice & Object class where used & Local quantum role & Finite-system status \\
\midrule
Underlying set of elements &
Bare finite sets and selected finite algebraic sets; in AQM-22, the field
\(\mathbb F_p\) viewed as a \(p\)-element set. &
Indexes all-map field labels such as
\(\operatorname{Map}(\mathbb F_p,\mathbb F_p^2)\), or the structureless
finite-set carriers of AQM-67--AQM-71 after extra data are supplied. &
Finite when the underlying set and target/carrier choices are finite.  In
AQM-22 this gives \(p\) qudits for the \(p\)-element set, not the one-point
spectrum. \\
\addlinespace
Rational points \(X(k)\) &
Schemes or varieties over a fixed field \(k\), especially finite-field
examples when a finite rational support is intentionally selected. &
Gives finite truncations such as the \(q^2\)-site rational-support subtheory
inside AQM-30.  It can also feed the finite-set field-label workflow after a
target \(V\), subspace \(E\), and pairing are chosen. &
Finite for the finite rational supports used in the report.  For an affine
finite-type presentation \(R=k[x_1,\ldots,x_n]/I\) with
\(k=\mathbb F_q\), the local finite-tuple derivation places \(X(k)\) inside
the finite set \(k^n\).  It is still only a truncation when non-rational
closed points also exist, as in AQM-30. \\
\addlinespace
Closed points &
Finite spectra with finite residue fields in AQM-23 and convention
\texttt{(u)}; affine finite-type \(\mathbb F_q\)-schemes in AQM-28--AQM-30
and convention \texttt{(y)}; closed prime points of \(\Spec\mathbb Z\) in
AQM-45 and convention \texttt{(ai)}. &
Each selected closed point \(x\) contributes the finite residue field
\(\kappa(x)\), the label space \(\kappa(x)^2\), and the finite local carrier
\(\ell^2(\kappa(x))\).  Finite closed supports quantize by tensor products of
the local matrix algebras. &
Finite for a finite closed support, and finite globally when the selected
closed-point set is finite, as in \(\Spec(\mathbb Z/6\mathbb Z)\).  Infinite
closed-point sets give algebraic quasi-local filtered colimits, not one
finite Hilbert space. \\
\addlinespace
Generic points &
Generic points in the infinite boundary examples: \((0)\in\Spec\mathbb
F_q[t]\), \((0)\in\Spec\mathbb F_q[x,y]\), and
\(\eta=(0)\in\Spec\mathbb Z\).  AQM-30 also records curve-generic points. &
They control topology, closures, divisor geometry, and open-set behaviour.
AQM-45 records an optional separate \(\mathbb Q\)-Weyl sector for
\(\Spec\mathbb Z\), but not as a finite qudit site. &
Not finite Weyl sites under the current closed-point conventions.  Generic
and curve-generic residue fields in AQM-29, AQM-30, and AQM-45 are infinite
or rational-function fields; analytic completions and adelic replacements are
gaps. \\
\addlinespace
All scheme points &
The full topological space \(|X|\) of an affine scheme, used for Zariski
opens and for stating the geometry in AQM-28--AQM-30 and AQM-45. &
Provides the ambient topology and inclusion order.  It is not automatically
the quantum-site set unless every selected point also has a recorded finite
or otherwise represented residue-sector convention. &
Finite only in special finite-spectrum cases already covered by the closed
point row.  For \(\Spec\mathbb F_q[t]\), \(\Spec\mathbb F_q[x,y]\), and
\(\Spec\mathbb Z\), all scheme points include non-closed generic-type points,
so this row is not a finite kinematical system by default. \\
\bottomrule
\end{tabular}
\end{center}

\subsection{Closed-Point Finite Weyl Review}

The finite Weyl-site convention already has three local forms.

First, AQM-23 and convention \texttt{(u)} treat a finite affine spectrum with
finite residue fields by assigning one local finite-field phase space to each
prime point.  For \(\Spec(\mathbb Z/6\mathbb Z)\), the two points have
residue fields \(\mathbb F_2\) and \(\mathbb F_3\).  The label group is the
orthogonal direct sum \(\mathbb F_2^2\oplus\mathbb F_3^2\), while the Hilbert
carrier is the tensor product \(\ell^2(\mathbb F_2)\otimes
\ell^2(\mathbb F_3)\).  AQM-23 proves that the observable algebra is the full
matrix algebra \(M_6(\mathbb C)\), not a Hilbert direct-sum algebra.

Second, AQM-28 and convention \texttt{(y)} extend this to affine schemes
\(X=\Spec R\) of finite type over \(\mathbb F_q\).  The finite Weyl sites are
the closed points \(x\in |X|_{\rm cl}\).  The local residue field
\(\kappa(x)\) is finite by the Nullstellensatz locator cited in AQM-28, so
\[
  E_x=\kappa(x)^2,\qquad
  H_x=\ell^2(\kappa(x)),\qquad
  A_x=\operatorname{End}_{\mathbb C}(H_x)
\]
are finite.  For an open \(U\), the label group is the finite-support direct
sum over \(U\cap |X|_{\rm cl}\), and the observable algebra is the filtered
colimit over finite closed supports.

Third, AQM-45 and convention \texttt{(ai)} use the same finite-residue rule
for the closed prime points \(x_\ell=(\ell)\) of \(\Spec\mathbb Z\), with
\(\kappa(x_\ell)=\mathbb F_\ell\).  The global closed-prime algebra is the
algebraic infinite tensor product over finite prime supports.  A Hilbert-space
realization needs an added representation choice, such as the
reference-vector incomplete tensor representation recorded in AQM-45.

\subsection{Finite Kinematical Boundary}

The point-selection rule gives a finite kinematical system only after both
the indexing family and the per-point quantum data are finite.

\begin{itemize}
\item A finite underlying set with finite carriers or finite symplectic target
data gives finite carriers after the extra data are named.  AQM-22 is the
prime-field example where \(\mathbb F_p\) as a set gives \(p\) qudit sites.
\item A finite rational-point support gives a finite truncation.  AQM-30 uses
the \(k\)-rational points of the affine plane as a finite \(q^2\)-site
subtheory, not as the full closed-point theory.
\item A finite closed support gives a finite tensor product of residue-field
qudits.  If the whole selected closed-point set is finite, as in AQM-23, the
global system is finite.
\item Infinite closed-point sets, as in AQM-29, AQM-30, and AQM-45, give
algebraic quasi-local systems through filtered finite supports.  They are
not single finite Hilbert spaces.
\item Generic points, curve-generic points, and other infinite-residue
points do not become finite Weyl sites without a separate representation
convention.
\end{itemize}

\subsection{Gaps Carried Forward}

This shard does not choose a generic-point representation convention.  The
generic residue fields \(k(t)\), \(k(x,y)\), curve-function fields, and
\(\mathbb Q\) are outside the finite closed-point Weyl rule.  AQM-45 records
one optional algebraic \(\mathbb Q\)-sector for \(\Spec\mathbb Z\), but that
does not settle local-field, adelic, \(C^*\)-completion, or analytic
Stone--von Neumann hypotheses for the other infinite-residue cases.

All later QSA shards that say "point" must therefore name the row above:
underlying element, rational point, closed point, generic point, or all scheme
point.  If the row is not the closed finite-residue row, the shard must also
state the extra carrier, residue, completion, support, or truncation choice
before making any finite Weyl or Hilbert-space claim.
